Let's take a look at the four rules more precisely. As we already know, they are used in order to mimic the real behavior of the DPLL algorithm, since its
direct implementation it's not so efficient.
\begin{enumerate}
    \item \textbf{UnitPropagate}. \\ 
    Mimics the generation of a new unit clause in DPLL. In practice, what the process really does is to extend the list $M$ with some new literals, 
    already derived, incorporating them inside the entire formula. Formally, this rule can be described as:
    $$ M \rightarrow Ml $$
    This occurs if $l$ is \textbf{undefined} in $M$, not in the list, and the CNF formula contains a clause $C \vee l$ satisfying $M \models \neg C$.
    \item \textbf{Decide}. \\ 
    Mimics the choice $p$ in DPLL, when no UnitPropagate is possible. \textbf{Decide} rule does not need any propagation, it just selects which
    boolean variable to introduce and its truth value.
    $$ M \rightarrow Ml^d $$
    This is possible if literal $l$ is \textbf{undefined} in $M$. \vspace{7pt}

    Note: $l^d$ symbolizes the choice done, it does not rely on any other mathematical expression.
    \item \textbf{Backtrack}. \\ 
    Mimics backtracking to the negation of the last decision in case a branch is unsatisfiable. Formally, it's defined as follows:
    $$ Ml^dN \rightarrow M \neg l $$
    where:
    \begin{itemize}
        \renewcommand{\labelitemi}{-}
        \item $l^d \rightarrow$ decision done
        \item $N \rightarrow$ set of derivations obtained
    \end{itemize} \vspace{3.5pt}

    This rule requires that $Ml^dN \models \neg C$ for a clause $C$ in the CNF formula and $N$ does not contain any literal coming from a decision. \vspace{7pt}

    If we reach a contradiction, we have to undo everything that came up from the last decision, deleting $N$ and negating $l$ (\textit{the truth value of the 
    literal $l$ does not matter, it can be either true or false}). \vspace{7pt}

    Note: \textbf{Backtrack} and \textbf{UnitPropogate} are little bit different. By UnitPropagate we can do something: incorporating literals $l$ to the formula we can still 
    satisfy the original clause; while in Backtrack there is no way to satisfy the clause $C$.
    \item \textbf{Fail}. \\ 
    Mimics the end of DPLL algorithm when every branch, and so the CNF formula, is unsatisfiable.
    $$ M \rightarrow \textit{fail} $$
    This happens if $M \models \neg C$ for a clause $C$ in the CNF formula and $M$ does not contain decision literals. \vspace{7pt}

    In practice, we have a clause that cannot be falsified and $M$ does not contain any decision $d$ left to backtrack.
\end{enumerate}